% ── Qualitative Examples (Appendix) ─────────────────────────────────────
% Requires: tcolorbox, xcolor, pifont packages
% \definecolor{darkgreen}{RGB}{0,128,0} in preamble

\section{Qualitative Examples}
\label{app:examples}

The following examples are drawn directly from model responses on \isosci{}
pairs and illustrate the three findings reported in Section~\ref{sec:results}.
Each box shows the source and target problems, the gold answers, and the
extracted model responses. \cmark{} = correct; \xmark{} = incorrect.

\begin{tcolorbox}[
  colback=gray!5, colframe=gray!40,
  title={\textbf{Example 1: Knowledge-Dependent Gain}},
  fonttitle=\small\bfseries, breakable]

\textbf{Finding:} \textit{o3-mini (R) solves the source problem but fails the isomorphic target — same reasoning structure, different domain knowledge.}
\\[4pt]
\textbf{Mapping:} Physics $\to$ Chemistry \quad
\textbf{Structure:} \texttt{unknown}

\vspace{6pt}
\textbf{Source problem (Physics):}\\
\begin{quote}\small Consider three identical, ideal capacitors. The first capacitor is charged to a voltage and then disconnected from the battery. The other two capacitors, initially uncharged and connected in series, are then connected across the first capacitor. What is the final voltage on the first capacitor?
A) V\\
\textit{Gold answer: D) 2V\_0/3}\end{quote}

\textbf{Target problem (Chemistry):}\\
\begin{quote}\small Consider three identical containers of ideal gas, each with volume V and initially at the same temperature T. The first container is filled with n moles of gas at pressure P\_0. The other two containers, initially evacuated and connected in series (sharing volume), are then connected to the first co\\
\textit{Gold answer: 2P\_0/3}\end{quote}

\vspace{4pt}
\textbf{o3-mini (R) on source} \textcolor{darkgreen}{\cmark} Extracted: \texttt{2V₀/3}
\\[-2pt]
\textbf{o3-mini (R) on target} \textcolor{red}{\xmark} Extracted: \texttt{P₀/3}
\begin{quote}\scriptsize\ttfamily [...] ombined system is given by the ideal gas law:
  P\_final = (nRT)/(3V).

• Substituting n from above gives:
  P\_final = [(P₀V)/(RT) × RT] / (3V) = P₀/3.

Since pressure equalizes throughout all connected containers, the final pressure in container 1 is P₀/3.

**Final Answer:** P₀/3\end{quote}
\\[-2pt]
\textbf{GPT-4o-mini (S) on source} \textcolor{red}{\xmark} Extracted: \texttt{The final voltage on the first capacitor is:}
\\[-2pt]
\textbf{GPT-4o-mini (S) on target} \textcolor{red}{\xmark} Extracted: \texttt{- Therefore, the final pressure in the first container after the gas has expande}
\\[-2pt]
\end{tcolorbox}
\vspace{8pt}
\begin{tcolorbox}[
  colback=gray!5, colframe=gray!40,
  title={\textbf{Example 2: Pure Knowledge Bottleneck}},
  fonttitle=\small\bfseries, breakable]

\textbf{Finding:} \textit{Both models solve the source problem correctly but fail the isomorphic target — neither possesses the target domain knowledge.}
\\[4pt]
\textbf{Mapping:} Physics $\to$ Chemistry \quad
\textbf{Structure:} \texttt{formula recall and substitute}

\vspace{6pt}
\textbf{Source problem (Physics):}\\
\begin{quote}\small Let \$X\$ equal the maximal oxygen intake of a human on a treadmill, where the measurements are in milliliters of oxygen per minute per kilogram of weight. Assume that, for a particular population, the mean of \$X\$ is \$\textbackslash\{\}mu=\$ 54.030 and the standard deviation is \$\textbackslash\{\\
\textit{Gold answer: 0.6247}\end{quote}

\textbf{Target problem (Chemistry):}\\
\begin{quote}\small Let Y equal the molar concentration of a sodium chloride solution prepared in a chemistry laboratory, where measurements are in moles per liter. Assume that, for a particular preparation method, the mean of Y is μ = 0.850 M and the standard deviation is σ = 0.042 M. Let Ȳ be the sample mean of a ran\\
\textit{Gold answer: 0.6247}\end{quote}

\vspace{4pt}
\textbf{o3-mini (R) on source} \textcolor{darkgreen}{\cmark} Extracted: \texttt{0.625 approximately.}
\\[-2pt]
\textbf{GPT-4o-mini (S) on source} \textcolor{darkgreen}{\cmark} Extracted: \texttt{0.6251}
\\[-2pt]
\textbf{o3-mini (R) on target} \textcolor{red}{\xmark} Extracted: \texttt{0.879 (approximately)}
\begin{quote}\scriptsize\ttfamily [...] 
Thus, the probability is given by
  P(–2 ≤ Z ≤ 1.2857) = P(Z ≤ 1.2857) – P(Z ≤ –2).

Using standard normal tables:
 • P(Z ≤ 1.29) is approximately 0.9015.
 • P(Z ≤ –2) is about 0.0228.

Subtracting these gives:
  0.9015 – 0.0228 ≈ 0.8787.

**Final Answer:** 0.879 (approximately)\end{quote}
\\[-2pt]
\textbf{GPT-4o-mini (S) on target} \textcolor{red}{\xmark} Extracted: \texttt{0.8787}
\begin{quote}\scriptsize\ttfamily [...] robability \textbackslash\{\}( P(0.836 \textbackslash\{\}leq \textbackslash\{\}bar\{Y\} \textbackslash\{\}leq 0.859) \textbackslash\{\}) is given by:
   \textbackslash\{\}[
   P(Z\_1 < Z < Z\_2) = P(Z < 1.29) - P(Z < -2.00)
   \textbackslash\{\}]
   Substituting the values:
   \textbackslash\{\}[
   P(0.836 \textbackslash\{\}leq \textbackslash\{\}bar\{Y\} \textbackslash\{\}leq 0.859) = 0.9015 - 0.0228 = 0.8787
   \textbackslash\{\}]

Thus, the final answer is:

**Final Answer:** 0.8787\end{quote}
\\[-2pt]
\end{tcolorbox}
\vspace{8pt}
\begin{tcolorbox}[
  colback=gray!5, colframe=gray!40,
  title={\textbf{Example 3: Reasoning Mode Helps (Source Domain Only)}},
  fonttitle=\small\bfseries, breakable]

\textbf{Finding:} \textit{o3-mini (R) solves the source problem correctly while GPT-4o-mini (S) fails — extended reasoning aids knowledge retrieval for the source domain formula.}
\\[4pt]
\textbf{Mapping:} Physics $\to$ Chemistry \quad
\textbf{Structure:} \texttt{unknown}

\vspace{6pt}
\textbf{Source problem (Physics):}\\
\begin{quote}\small Consider three identical, ideal capacitors. The first capacitor is charged to a voltage and then disconnected from the battery. The other two capacitors, initially uncharged and connected in series, are then connected across the first capacitor. What is the final voltage on the first capacitor?
A) V\\
\textit{Gold answer: D) 2V\_0/3}\end{quote}

\textbf{Target problem (Chemistry):}\\
\begin{quote}\small Consider three identical containers of ideal gas, each with volume V and initially at the same temperature T. The first container is filled with n moles of gas at pressure P\_0. The other two containers, initially evacuated and connected in series (sharing volume), are then connected to the first co\\
\textit{Gold answer: 2P\_0/3}\end{quote}

\vspace{4pt}
\textbf{o3-mini (R)} \textcolor{darkgreen}{\cmark} Extracted: \texttt{2V₀/3}
\begin{quote}\scriptsize\ttfamily [...] final.
By charge conservation, the total charge remains the same:
  Q\_initial = Q₁ + Q\_series  ⇒  C·V₀ = C·V\_final + (C/2)·V\_final = (3C/2)·V\_final.
Solving for V\_final gives:
  V\_final = (2/3)·V₀.

Thus, the final voltage on the first capacitor is 2V₀/3.

**Final Answer:** 2V₀/3\end{quote}
\\[-2pt]
\textbf{GPT-4o-mini (S)} \textcolor{red}{\xmark} Extracted: \texttt{The final voltage on the first capacitor is:}
\begin{quote}\scriptsize\ttfamily [...] ):
   \textbackslash\{\}[
   V\_0 = 2 V\_f
   \textbackslash\{\}]
   Therefore, we can solve for \textbackslash\{\}( V\_f \textbackslash\{\}):
   \textbackslash\{\}[
   V\_f = \textbackslash\{\}frac\{V\_0\}\{2\}
   \textbackslash\{\}]

7. **Final Answer**: The final voltage on the first capacitor is:
   \textbackslash\{\}[
   \textbackslash\{\}boxed\{\textbackslash\{\}frac\{V\_0\}\{2\}\}
   \textbackslash\{\}]

Thus, the final answer is:

**Final Answer:** C) \textbackslash\{\}( \textbackslash\{\}frac\{V\_0\}\{2\} \textbackslash\{\})\end{quote}
\\[-2pt]
\end{tcolorbox}
\vspace{8pt}
\begin{tcolorbox}[
  colback=gray!5, colframe=gray!40,
  title={\textbf{Example 4: Knowledge-Dependent Gain}},
  fonttitle=\small\bfseries, breakable]

\textbf{Finding:} \textit{Qwen3-32B thinking=ON (R) solves the source problem but fails the isomorphic target — same reasoning structure, different domain knowledge.}
\\[4pt]
\textbf{Mapping:} Physics $\to$ Chemistry \quad
\textbf{Structure:} \texttt{unknown}

\vspace{6pt}
\textbf{Source problem (Physics):}\\
\begin{quote}\small A liquid flows at a constant flow rate through a pipe with circular cross-sections of varying diameters. At one point in the pipe, the diameter is 2 cm and the flow speed is 18 m/s. What is the flow speed at another point in this pipe, where the diameter is 3 cm?
A) 4 m/s
B) 6 m/s
C) 8 m/s
D) 12 m/s\\
\textit{Gold answer: C) 8 m/s}\end{quote}

\textbf{Target problem (Chemistry):}\\
\begin{quote}\small A gas diffuses through a porous membrane at a constant molar flow rate. At one location in the membrane, the cross-sectional area is 4.0 cm² and the diffusion flux is 0.12 mol/(cm²·s). What is the diffusion flux at another location where the cross-sectional area is 6.0 cm²?
A) 0.04 mol/(cm²·s)
B) 0.\\
\textit{Gold answer: C) 0.08 mol/(cm²·s)}\end{quote}

\vspace{4pt}
\textbf{Qwen3-32B thinking=ON (R) on source} \textcolor{darkgreen}{\cmark} Extracted: \texttt{C) 8 m/s}
\\[-2pt]
\textbf{Qwen3-32B thinking=ON (R) on target} \textcolor{red}{\xmark} Extracted: \texttt{\$\$}
\begin{quote}\scriptsize\ttfamily [...] frac\{4.0\}\{6.0\} = 0.12 \textbackslash\{\}cdot \textbackslash\{\}frac\{2\}\{3\} = 0.08 \textbackslash\{\}, \textbackslash\{\}text\{mol/(cm\}\textasciicircum{}2 \textbackslash\{\}cdot \textbackslash\{\}text\{s)\}
\$\$

This result makes sense intuitively: as the cross-sectional area increases, the same molar flow rate is spread over a larger area, so the flux decreases.

---

\#\#\# Final Answer

\$\$
\textbackslash\{\}boxed\{C\}
\$\$\end{quote}
\\[-2pt]
\textbf{Qwen3-32B thinking=OFF (S) on source} \textcolor{red}{\xmark} Extracted: \texttt{\$\$}
\\[-2pt]
\textbf{Qwen3-32B thinking=OFF (S) on target} \textcolor{darkgreen}{\cmark} Extracted: \texttt{C) 0.08 mol/(cm²·s)}
\\[-2pt]
\end{tcolorbox}
\vspace{8pt}
\begin{tcolorbox}[
  colback=gray!5, colframe=gray!40,
  title={\textbf{Example 5: Pure Knowledge Bottleneck}},
  fonttitle=\small\bfseries, breakable]

\textbf{Finding:} \textit{Both models solve the source problem correctly but fail the isomorphic target — neither possesses the target domain knowledge.}
\\[4pt]
\textbf{Mapping:} Physics $\to$ Biology \quad
\textbf{Structure:} \texttt{unknown}

\vspace{6pt}
\textbf{Source problem (Physics):}\\
\begin{quote}\small Planetary rings are
A) known to exist for all of the jovian planets.
B) composed of a large number of individual particles that orbit their planet in accord with Kepler's third law.
C) nearer to their planet than any of the planet's large moons.
D) all of the above\\
\textit{Gold answer: D) all of the above}\end{quote}

\textbf{Target problem (Biology):}\\
\begin{quote}\small Homologous chromosomes are
A) present as pairs in diploid organisms, with one chromosome inherited from each parent.
B) similar in size, shape, and genetic content, carrying genes for the same traits at corresponding loci.
C) separated during meiosis I, resulting in haploid daughter cells with one c\\
\textit{Gold answer: D) all of the above}\end{quote}

\vspace{4pt}
\textbf{Qwen3-32B thinking=ON (R) on source} \textcolor{darkgreen}{\cmark} Extracted: \texttt{D) all of the above}
\\[-2pt]
\textbf{Qwen3-32B thinking=OFF (S) on source} \textcolor{darkgreen}{\cmark} Extracted: \texttt{D) all of the above}
\\[-2pt]
\textbf{Qwen3-32B thinking=ON (R) on target} \textcolor{red}{\xmark} Extracted: \texttt{D) all of the above.}
\begin{quote}\scriptsize\ttfamily [...] s. This process reduces the chromosome number by half, resulting in haploid cells, which is essential for sexual reproduction.

---

Since all three statements (A, B, and C) accurately describe homologous chromosomes, the correct choice is:

**Final Answer:** D) all of the above.\end{quote}
\\[-2pt]
\textbf{Qwen3-32B thinking=OFF (S) on target} \textcolor{red}{\xmark} Extracted: \texttt{D) all of the above.}
\begin{quote}\scriptsize\ttfamily [...] n separated into different daughter cells. This process is crucial for producing gametes with a haploid number of chromosomes.

---

Since all three statements (A, B, and C) accurately describe homologous chromosomes, the correct choice is:

**Final Answer:** D) all of the above.\end{quote}
\\[-2pt]
\end{tcolorbox}
\vspace{8pt}
\begin{tcolorbox}[
  colback=gray!5, colframe=gray!40,
  title={\textbf{Example 6: Reasoning Mode Helps (Source Domain Only)}},
  fonttitle=\small\bfseries, breakable]

\textbf{Finding:} \textit{Qwen3-32B thinking=ON (R) solves the source problem correctly while Qwen3-32B thinking=OFF (S) fails — extended reasoning aids knowledge retrieval for the source domain formula.}
\\[4pt]
\textbf{Mapping:} Physics $\to$ Chemistry \quad
\textbf{Structure:} \texttt{unknown}

\vspace{6pt}
\textbf{Source problem (Physics):}\\
\begin{quote}\small A liquid flows at a constant flow rate through a pipe with circular cross-sections of varying diameters. At one point in the pipe, the diameter is 2 cm and the flow speed is 18 m/s. What is the flow speed at another point in this pipe, where the diameter is 3 cm?
A) 4 m/s
B) 6 m/s
C) 8 m/s
D) 12 m/s\\
\textit{Gold answer: C) 8 m/s}\end{quote}

\textbf{Target problem (Chemistry):}\\
\begin{quote}\small A gas diffuses through a porous membrane at a constant molar flow rate. At one location in the membrane, the cross-sectional area is 4.0 cm² and the diffusion flux is 0.12 mol/(cm²·s). What is the diffusion flux at another location where the cross-sectional area is 6.0 cm²?
A) 0.04 mol/(cm²·s)
B) 0.\\
\textit{Gold answer: C) 0.08 mol/(cm²·s)}\end{quote}

\vspace{4pt}
\textbf{Qwen3-32B thinking=ON (R)} \textcolor{darkgreen}{\cmark} Extracted: \texttt{C) 8 m/s}
\begin{quote}\scriptsize\ttfamily [...] ght) = \textbackslash\{\}frac\{72\}\{9\} = 8 \textbackslash\{\}, \textbackslash\{\}text\{m/s\}
\$\$

---

\#\#\# Step 3: Interpret the Result

The flow speed decreases at the wider section of the pipe, which is consistent with the principle of continuity. The correct answer is:

\$\$
\textbackslash\{\}boxed\{8 \textbackslash\{\}, \textbackslash\{\}text\{m/s\}\}
\$\$

---

**Final Answer:** C) 8 m/s\end{quote}
\\[-2pt]
\textbf{Qwen3-32B thinking=OFF (S)} \textcolor{red}{\xmark} Extracted: \texttt{\$\$}
\begin{quote}\scriptsize\ttfamily [...] s. Since the flow rate must remain constant, the flow speed must **decrease**. The calculated speed of 8 m/s is indeed less than the initial speed of 18 m/s, which is consistent with the physical behavior of incompressible fluids in pipes.

---

\#\#\# Final Answer:

\$\$
\textbackslash\{\}boxed\{C\}
\$\$\end{quote}
\\[-2pt]
\end{tcolorbox}
\vspace{8pt}
\begin{tcolorbox}[
  colback=gray!5, colframe=gray!40,
  title={\textbf{Example 7: Structure-Invariant Gain (Rare: 4.8\% of gains)}},
  fonttitle=\small\bfseries, breakable]

\textbf{Finding:} \textit{Qwen3-32B thinking=ON (R) solves both source and target correctly while Qwen3-32B thinking=OFF (S) fails both — one of the few cases where reasoning improvement transfers across domain knowledge.}
\\[4pt]
\textbf{Mapping:} Physics $\to$ Earth Science \quad
\textbf{Structure:} \texttt{unknown}

\vspace{6pt}
\textbf{Source problem (Physics):}\\
\begin{quote}\small A table supports a wooden block placed on the tabletop. Which fundamental force of nature is responsible for this interaction, and why?
A) The electric force, because the protons in the nuclei of the top atomic layer of the table repel the nuclei in the bottom atomic layer of the wood.
B) The gravit\\
\textit{Gold answer: C) The electric force, because the outer electrons in the top atomic layer of the table repel the outer electrons in the bottom atomic layer of the wood.}\end{quote}

\textbf{Target problem (Earth Science):}\\
\begin{quote}\small Water molecules adhere to the surface of soil particles in the vadose zone above the water table. Which fundamental force of nature is primarily responsible for this adhesion, and why?
A) The gravitational force, because the mass of water molecules is attracted to the larger mass of soil particles b\\
\textit{Gold answer: B) The electromagnetic force, because polar water molecules are attracted to charged or polar sites on soil particle surfaces through electrostatic interactions.}\end{quote}

\vspace{4pt}
\textbf{Qwen3-32B thinking=ON (R) on source} \textcolor{darkgreen}{\cmark} Extracted: \texttt{C) The electric force, because the outer electrons in the top atomic layer of th}
\\[-2pt]
\textbf{Qwen3-32B thinking=ON (R) on target} \textcolor{darkgreen}{\cmark} Extracted: \texttt{B) The electromagnetic force, because polar water molecules are attracted to cha}
\\[-2pt]
\textbf{Qwen3-32B thinking=OFF (S) on source} \textcolor{red}{\xmark} Extracted: \texttt{\textbackslash\{\}text\{C}
\\[-2pt]
\textbf{Qwen3-32B thinking=OFF (S) on target} \textcolor{red}{\xmark} Extracted: \texttt{B}
\\[-2pt]
\end{tcolorbox}
% ─────────────────────────────────────────────────────────────────────────
